% !TEX root = ../nc-quadric-surface.tex
\section{重み付き四周期表示の偶数環}
\label{sec:weighted-inverse}

以下、$D$ を局所有限で正に次数付けされた連結 $\mathbb Z$-algebra とし、
$P_j=e_jD$ と書く。$D$ は左右 noetherian とし、右単純加群 $S_j$ は
\[
 0\longrightarrow P_{j+4}\longrightarrow
 P_{j+3}^2\oplus P_{j+2}^{d_j}\longrightarrow
 P_{j+1}^2\oplus P_{j+2}^{d_j}\longrightarrow
 P_j\longrightarrow S_j\longrightarrow0
 \tag{*}\label{eq:weighted-inverse-rank-resolution}
\]
という極小分解を持つとする。ここで $d_j=0$（$j$ 奇数）、
$d_j\in\{0,1\}$（$j$ 偶数）、$d_{j+2}=d_j$ とする。
この節ではこの完全性を仮定して逆向きに何が従うかを調べる。

\begin{proposition}[周期偶数環の Koszul 性]
\label{prop:weighted-even-koszul}
$D$ の偶数部分を二添字の周期移動で次数付けして得られる環を $B$ とする。
すると $B$ は左右 noetherian、一次生成、二次 Koszul であり、
$H_B(z)=(1+z)/(1-z)^3$ である。
\end{proposition}

\begin{proof}
極小分解の第一項から、$D$ は隣接する二つの生成元と
偶数添字から始まる距離二の生成元で生成される。
偶数添字間の道は長さ一か二で次の偶数添字を通るから、
$B$ は一次生成である。偶数部分の斉次右イデアルを
$P_j$（$j$ 偶数）の部分加群として $D$ へ拡張し、
偶数成分に制限すると元のイデアルに戻る。
実際、拡張で最終添字が偶数になる項に使われる積は
両端とも偶数の成分だけである。よって斉次右イデアルの昇鎖は停止する。
左側も同様であり、一般のイデアルには最高次の初期項を用いる。

$h_{jk}=\dim D_{jk}$ と書く。極小分解の各成分で Euler 標数を取り、
同じ添字に現れる距離二の項を消去すると、
\[
 h_{jk}-2h_{j+1,k}+2h_{j+3,k}-h_{j+4,k}=\delta_{jk}.
\]
$h_{jk}=0$（$j>k$）から順に解けば、$h_{jk}$ は $k-j$ のみの関数となり、
その母関数は $(1-2z+2z^3-z^4)^{-1}$ である。
偶数差 $2n$ の係数は $(n+1)^2$ になる。

奇数 $i$ に対し、$N_i=P_i|_{\mathrm{even}}$ と置く。
式 \eqref{eq:weighted-inverse-rank-resolution} を偶数成分に制限すると
$S_i$ は零になるため、
\[
 0\to N_{i+4}\to (P_{i+3}|_{\mathrm{even}})^2
   \to(P_{i+1}|_{\mathrm{even}})^2\to N_i\to0
\]
を得る。$i$ を四ずつ増やしてこの列をつなぐと、
$N_i$ の自由分解の各項は、最初の生成次数から一つずつ増加する次数で生成される。
特に $N_1$ の左端の生成次数は $1$、$N_3$ では $2$ である
（偶数添字 $2j$ を次数 $j$ と数える）。

次に $j=0$ の式 \eqref{eq:weighted-inverse-rank-resolution} を
偶数成分に制限し、$k$ の分解を得る。
複体の第 $p=0,1,2,3$ 項は順に
$B$、$N_1^2\oplus B(-1)^{d_0}$、
$N_3^2\oplus B(-1)^{d_0}$、$B(-2)$ である。
この複体に $-\otimes_B^{\mathbf L} k$ を適用した
スペクトル系列の第一頁では、各第 $p$ 項の
$\operatorname{Tor}_q^B(-,k)$ は、上で得た線形分解により
内部次数が $p+q$ 以下にある。
従って $\operatorname{Tor}_n^B(k,k)$ は内部次数 $n$ 以下にある。
一方、$B$ は連結一次生成なので、極小分解から同じ Tor は内部次数 $n$ 以上にある。
よって $B$ は Koszul であり、特に二次である。
\end{proof}

\begin{proposition}[残る許容条件への還元]
\label{prop:weighted-inverse-reduction}
命題~\ref{prop:weighted-even-koszul} の仮定に加え、
$D$ を四周期、次元三・parameter 四の AS 正則環とする。
偶数部分の次数付けを与える周期移動の二乗は、
$D$ の逆 Nakayama 同型を偶数部分に制限したものと
対角型自己同型を除いて一致するとする。
このとき、次の条件が成立すれば $(\qgr D,\pi e_0D)\in\HH$ である。
\begin{enumerate}
\item $B$ は domain である。
\item ある正次数の非零斉次元 $f\in B$ に対して
$\pi(B/fB)$ は $\qgr B$ で有限長でない。
\item 定理~\ref{thm:finite-ambient-test} の有限階数条件を満たす
超平面がある。
\end{enumerate}
これらは、この周期移動を適合根とする逆構成に必要でもある。
\end{proposition}

\begin{proof}
偶数部分への制限により $\qgr D\simeq\qgr B$ となる。
ここでは制限は exact であり、左随伴は偶数頂点の表現可能加群を
対応する $P_{2j}$ に送る加群の延長である。
この随伴の単位は偶数成分で同型であり、
余単位の核と余核は偶数成分が零となる。
偶数成分が零の有限生成加群は、奇数始点の任意の正距離の道が
最初に偶数頂点を通るため、有限次元である。
また有限生成加群の偶数部分も有限生成であり、
奇数頂点に置かれた生成元については、次の偶数頂点への二本の矢でその制限を生成できる。
このためこれらの操作は tails において互いに逆である。
添字零の表現可能加群は $\pi B$ に移る。

\cite[Theorems 4.19, 6.8 and 6.11, Proposition 4.11]{MN2025} から、
$D$ 側の標識は exceptional で、四添字の Nakayama 移動が
反 Serre 関手となる。仮定した対角型の差は加群の各成分を定数倍すれば
自然同型になるので、偶数部分のシフトはその平方根である。
$\qgr D$ の豊富な表現可能加群列は、奇数項を次の偶数項から
生成できるため、偶数部分列でも豊富である。
表現可能加群間の Hom が $D$ の成分と一致することから、
この根の切断環は全次数で $B$ となる。
従って条件 (3) が所要の正則正規被覆を与える。

条件 (1) を仮定する。条件 (3) の被覆からすべての有限生成加群の
Hilbert 多項式が得られ、次数が高々一の対象は $\cC_1$ に入る
（命題~\ref{prop:sklyanin-h2} と同じ長さの議論）。
$f$ は正規である必要はなく、
$f^jB/f^{j+1}B\simeq(B/fB)(-j\deg f)$ が domain 性から従う。
条件 (2) のとき、この列は $\cC/\cC_0$ でも真に下降するから
$\pi B\notin\cC_1$ である。
$\pi B$ の任意の非零部分対象が $\pi B$ のシフトを含むことと、
それによる余商の Hilbert 多項式の次数が高々一であることから、
rank one torsion-free 性が従う。

逆に、適合根を得たなら $B$ の domain 性と被覆の有限条件は既に示した。
もしすべての非零斉次 $f$ に対して $\pi(B/fB)$ が有限長なら、
$\pi B$ の任意の非零部分対象による商は有限長となる。
$\cC/\cC_0$ の部分対象を持ち上げると、
そこにおける $\pi B$ の像は零または単純となり、
$\pi B\in\cC_1$ となって矛盾する。
\end{proof}

この還元では $D$ の四周期最小分解の完全性と noetherian 性を仮定した。
残る分類課題は、前節の距離二・距離三の関係からこれらを保証する条件、
および周期偶数環の domain 性と条件 (2) の具体的な判定である。
